\IEEEPARstart{A}{utonomous} vehicles plan around forecasts of what nearby road users will do. A forecast is only as
useful as the statement it makes about its own reliability: a planner that knows a predicted region contains the true
future with probability at least $1-\alpha$ can keep clear of that region and of nothing larger. Conformal prediction
(CP) supplies such a statement for any black-box predictor, with a finite-sample guarantee that requires nothing beyond
exchangeability of calibration and test data \cite{vovk2005algorithmic,angelopoulos2023gentle}, and it has been adopted
for trajectory prediction and safe planning
\cite{lindemann2023safe,sun2024copula,huang2025cuqds}.

Exchangeability is exactly what deployment breaks. A forecaster is calibrated on the data at hand---one set of cities,
one sensor suite, one annotation pipeline---and then driven somewhere else. That accuracy drops when a trajectory
predictor moves between datasets is documented \cite{feng2024unitraj}, and that predictive uncertainty degrades under
dataset shift is well known for classifiers \cite{ovadia2019can}. What is not known is what happens to the
\emph{guarantee itself} when a conformally calibrated forecaster is moved, without any new labels, to a different real
driving dataset: whether the stated $1-\alpha$ survives, how badly it fails, and whether the standard remedies for shift
repair it.

The conformal literature offers several remedies, each resting on an assumption. Weighted CP restores validity under
covariate shift when the likelihood ratio is known or can be estimated \cite{tibshirani2019covariate}; more general
results bound the coverage loss by a distance between calibration and test distributions
\cite{barber2023beyond}; online methods track shift by feeding realised errors back into the threshold
\cite{gibbs2021adaptive,angelopoulos2023pid}, an idea carried to trajectory prediction \cite{huang2025cuqds} and to
robotic settings more broadly \cite{tumu2026adaptnc}; and a generative-prior approach has been proposed for planning
under environment shift \cite{rahaman2026shift}. What these do not establish is the \emph{static, label-free,
dataset-to-dataset} setting in which a forecaster is frozen, calibrated once, and deployed: how large the failure is on
real data, whether the covariate-shift remedy suffices, and how many labelled target scenes it takes to detect or repair
the failure.

\paragraph*{Contributions}
(i) We give the first measurement of zero-label conformal coverage transfer between real driving datasets, in both
directions between Argoverse~2 and nuScenes and, pre-registered before it was run, on a third dataset, Waymo Open
Motion. The gap reaches \gapRawConf\ points at a nominal 90\% level on a model trained to match published
state-of-the-art accuracy, and up to \wmGapAvW\ points on Waymo; it grows, not shrinks, as the point predictor
improves. (ii) We diagnose the failure: covariate reweighting does not repair it and makes it worse on every model
tested, and we give a non-identifiability result (Proposition~\ref{prop:ident}) explaining why a label-free method is
limited in principle. Controlled input perturbations show that measured differences between the datasets explain
only a minority of the gap, and localise most of it to turning manoeuvres and left-hand traffic. (iii) We give a
repair with a finite-sample, training-conditional guarantee (Propositions~\ref{prop:cor}
and~\ref{prop:scene}) that holds when labels arrive scene by scene, the realistic case for driving data, where the
natural per-agent guarantee is shown empirically to fail. With 30--40 labelled target scenes it restores coverage
close to nominal at a moderate cost in region size, confirmed on gate-passing models across both transfer directions
and all three datasets. (iv) We release the code, the pre-registered hypotheses, and an append-only experiment log
recording every result, including the negative and partially-supported ones.
